\noindent \textit{\textcolor{red!60!black}{
1. Framing:\\
The paper is closely related to firm-level uncertainty after Brexit studied by Hassan et al (2019, not cited). While a whole cottage industry has followed the firm-level approach pioneered by Hassan et al, I am not aware of a paper that focuses on an uncertainty shock to financial arrangements. Therefore, identifying a case of firm-level uncertainty about their financial contracts is a significant contribution to the literature. While the authors briefly point to Bloom (2009), the authors should highlight these firm-level differences in uncertainty much more and emphasize that these differences matter for aggregate outcomes during the Great Depression. Given the nature of the shock and the resulting litigation, I don't think the effects are due to 'leverage risk' but to uncertainty about leverage. By reframing the paper and linking it to the literature on firm-level uncertainty, the paper would have a larger impact.
}}
 \bigskip

\noindent Thank you for this insightful comment, which has substantially shaped our revised manuscript. Following your suggestion, we have engaged more deeply with the literature on policy uncertainty and its effects on firm outcomes \citep{Bloom2009, Baker2016, Hassan2019}. In particular, we read \cite{Hassan2019} carefully, along with the broader literature that has followed their work. This engagement has clarified how our contribution fits within this active area of research.

\cite{Hassan2019} pioneers a firm-level approach to measuring political uncertainty, using textual analysis of earnings calls to identify heterogeneous exposure to policy shocks across firms. A key insight from this work is that aggregate measures of uncertainty can mask substantial firm-level variation in exposure, and that this variation has real consequences for firm decisions. We now recognize that our setting offers a complementary approach to measuring firm-level exposure: rather than inferring exposure from managerial discussions, we measure it through the gold content of firms' liabilities. This parallel is now explicit in the literature review, where we write that ``in the same spirit as \cite{Hassan2019}, who measure firm-level exposure to policy shocks using textual analysis of earnings call transcripts, we identify firm-level exposure to legal uncertainty using the gold content of firms' liabilities.''

More broadly, the influence of your suggestion can be seen throughout the paper, starting with its new title: ``Debt Overhang During Policy Uncertainty: The Case of Gold Clauses in the 1930s.'' In the introduction, we now position our study as ``closest to the literature on firm-level political uncertainty pioneered by \cite{Hassan2019}'' and emphasize that the abrogation represents ``political intervention in private debt contracts.'' While \cite{Hassan2019} documents the real effects of uncertainty in modern episodes such as Brexit, U.S. elections, and Italian constitutional reforms, our study examines how similar mechanisms shaped firm decisions during the Great Depression---a pivotal historical episode that complements the modern evidence.

As you note in the second paragraph of your report, by studying the effects of debt overhang, we ``follow a well-trodden path.'' However, the unique features of our empirical setting also enable us to study how policy uncertainty that \textit{directly} impacts the right-hand side of the firm's balance sheet (through leverage uncertainty) also \textit{indirectly} impacts the left-hand side (through firms' investment decisions). While we interpret our results through the lens of debt overhang, in this episode it takes a unique form shaped by political decisions affecting firm-level choices. We now state in the introduction that ``we contribute to this literature by showing how an uncertainty shock to future liabilities delayed the recovery of investment during the Great Depression.''

